\chapter{Introduction}
\label{cha:introduction}

\section{Background and research scope}
\label{sec:intro_background}

% opening paragraph
The isolation of graphene demonstrated that reducing a material to the atomic scale can generate electronic, optical, and quantum properties that differ substantially from those of its bulk form~\cite{novoselov2004electric,geim2007rise}.
Since then, two-dimensional materials have become a major platform in condensed matter physics and materials science.
Reduced dimensionality reshapes the electronic states, enhances the role of surfaces and interfaces, and strengthens the connection between symmetry, bonding, and physical response.
These features allow two-dimensional systems to host properties that are difficult to obtain in conventional three-dimensional crystals.

Beyond graphene, the family of atomically thin materials has expanded to include elemental sheets commonly referred to as Xenes, as well as related light-element compound two-dimensional materials~\cite{butler2013progress,molle2018silicene}.
Light-element Xenes and related compounds provide a useful setting for studying the connection between bonding and physical response.
Their low atomic masses and diverse bonding environments give rise to structures with different electronic properties.
Depending on atomic arrangement and chemical modification, these systems may exhibit metallicity, semiconductivity, magnetism, superconductivity, non-trivial topology, and anisotropic optical response.
The relevant question is therefore how a particular structural change modifies these properties, and which calculations can establish that connection.

Boron-based two-dimensional materials occupy an important position in this landscape.
Unlike carbon, boron is electron deficient, and this electron deficiency favours multicentre bonding and structural diversity.
As a result, boron can form many competing atomic networks with distinct electronic and collective properties~\cite{sun2017two,han2023superconducting}.
Boron allotropes and boron-containing compounds therefore provide natural systems for examining how bonding topology controls physical response.
Beryllium-based two-dimensional materials offer a complementary setting.
Beryllium is the lightest alkaline-earth metal, and both experimental and theoretical studies have extended its investigation to atomically thin structures~\cite{ono2020dynamical,chahal2023beryllene,li2023coexistence}.
These studies motivate a comparison of how layer geometry and hydrogen functionalization modify the properties of an elemental light-metal system.

A central challenge in two-dimensional materials research is to tune intrinsic properties for specific physical functions.
Van der Waals heterostructure formation provides one route.
In such systems, different monolayers are stacked through weak interlayer interactions, and the interface can modify electronic states, charge transfer, optical transitions, and collective excitations while preserving important features of the individual layers~\cite{novoselov20162d}.
Surface functionalization provides another route.
Hydrogen adsorption or chemical termination can alter bonding, redistribute electronic states, and change the dielectric response.
We therefore use these strategies to examine how structural modification changes electronic structure and optical response.
For selected systems, we also investigate magnetism, superconductivity, and band topology.

\section{Literature review}
\label{sec:intro_lit_review}

% opening paragraph
The previous section introduced the materials and the structural changes considered in this thesis.
We now examine what earlier studies have established about these systems.
We first consider boron-based materials and graphene heterostructures, then turn to beryllium, and finally compare the computational approaches.
Throughout this review, we distinguish experimental observations from theoretical predictions and identify the questions that remain open.

\subsection{Boron-based materials and graphene heterostructures}
\label{subsec:lit_boron}

Firstly, the structure of a boron sheet is not specified by its elemental composition alone.
The arrangement of triangular units, vacancies, and connections between layers changes the bonding network.
Experimental work on Ag(111) has identified different borophene phases under different growth conditions~\cite{zhong2017metastable}.
Chen et al. subsequently reported the synthesis of bilayer borophene on Cu(111), extending the experimentally accessible structures beyond a single atomic layer~\cite{chen2022synthesis}.
These results show why the phase, layer number, and supporting surface need to be specified together.
A substrate-supported structure establishes an experimental material, but its bonding and electronic response need not be identical to those of an isolated sheet.
Therefore, a comparison with a freestanding calculation requires attention to both the atomic structure and its environment.

Theoretical studies have also examined how different boron networks support different electronic states.
Gao et al. calculated electron--phonon coupling in metallic $\beta_{12}$- and $\chi_3$-borophene and predicted superconductivity using the McMillan--Allen--Dynes approach~\cite{gao2017prediction}.
Zhu et al. used an evolutionary structure search to identify magnetic borophene structures~\cite{zhu2019magnetic}.
These investigations concern specified geometries and specified electronic states.
They therefore provide examples of the sensitivity of boron properties to the bonding network, rather than a prediction that all boron sheets share the same magnetic or superconducting behaviour.
For a different structural framework, the electronic configuration, phonons, and electron--phonon coupling must be examined again.

Next, incorporating carbon introduces another way to change the bonding and electronic structure.
Experimental studies have investigated epitaxial \ce{BC3} honeycomb sheets and their phonon properties~\cite{yanagisawa2004phonon,tanaka2005novel,yanagisawa2006phonon}.
Theoretical work has proposed additional two-dimensional boron carbides, including a semiconducting \ce{B4C3} monolayer~\cite{tian2019}.
These two types of evidence should be distinguished.
An experimental \ce{BC3} sheet on a substrate is a relevant precedent, while the properties of an ideal freestanding model remain predictions for that model.
Similarly, the predicted stability and band structure of \ce{B4C3} motivate further study but do not establish its experimental preparation.
For the present work, the selected \ce{BC3} and \ce{B4C3} structures provide semiconducting constituents that can be compared with metallic borophene and semimetallic graphene.

The interface introduces a further degree of freedom.
Liu and Hersam reported lateral and vertical integration of borophene with graphene, while Hou et al. investigated the preparation and humidity-sensing response of borophene--graphene heterostructures~\cite{liu2019borophene,hou2021borophene}.
These studies establish that borophene--graphene combinations are experimentally relevant.
However, their structures, environments, and measured properties do not answer every question about the ideal bilayers considered here.
In particular, comparisons between different boron-based constituents require the reference monolayers and their interfaces to be treated within a consistent computational setting.
We therefore compare graphene--borophene, graphene--\ce{BC3}, and graphene--\ce{B4C3} in chapter~\ref{cha:project_bilayers}.
The central question is how interfacial coupling modifies the electronic structure and optical response relative to the isolated constituents.
The contribution is a comparative mapping of these structure--property relations.

We now turn from interfaces between different sheets to dimensional reduction within a single boron framework.
Han et al. predicted a bulk orthorhombic allotrope, o-\ce{B14}, using first-principles calculations and evolutionary structure searches~\cite{han2023superconducting}.
Their search examined potassium--boron compounds, and removal of potassium from the predicted frameworks generated candidate boron allotropes.
The o-\ce{B14} structure contains pentagonal bipyramidal units and was predicted to support phonon-mediated superconductivity.
This is a theoretical bulk reference, not an experimentally synthesized starting material.
It also differs structurally from the planar borophenes discussed above.
Reducing this framework to one or two layers changes the coordination of the surface atoms, while hydrogen termination modifies their bonding further.
The bulk calculation therefore leaves open how these changes affect the electronic ground state and physical response of the resulting sheets.

Hydrogenation studies on other boron structures show why the outcome cannot be inferred from hydrogen addition alone.
Hou et al. reported crystalline semiconducting hydrogenated borophene produced through a chemical preparation route~\cite{hou2020ultrastable}.
Li et al. instead exposed borophene to atomic hydrogen under ultrahigh vacuum and obtained borophane polymorphs that retained metallicity~\cite{li2021synthesis}.
The structures and preparation conditions differ, so these results do not constitute conflicting measurements of the same material.
They show that the parent network and hydrogen arrangement are part of the physical problem.
Neither result directly establishes the behaviour of hydrogen-terminated o-\ce{B14}.
We therefore examine dimensional reduction and hydrogen termination within this particular framework in chapter~\ref{cha:project_boron}, using the bulk structure and the corresponding pristine sheets as references.
The comparison concerns electronic and magnetic states, electron--phonon properties, and optical response.

Figure~\ref{fig:intro} shows the hydrogen-terminated bilayer model considered in this thesis.

\begin{figure}[!htbp]
\centering
    \includegraphics[width=0.6\linewidth]{figures_ch1/intro.png}
\caption{Atomic structure of hydrogen-terminated bilayer o-\ce{B14}. Large coloured spheres represent boron atoms, with different colours denoting crystallographically inequivalent boron sites. Small light-purple spheres represent hydrogen atoms. The figure illustrates a computational structural model considered in chapter~\ref{cha:project_boron}.}
\label{fig:intro}
\end{figure}

\FloatBarrier
\subsection{Beryllene and hydrogen functionalization}
\label{subsec:lit_beryllene}

The previous subsection considered how interfaces and dimensional reduction modify boron-based systems.
We now extend the discussion to beryllium, where the choice of layer geometry is itself an important part of the investigation.
Ono compared the dynamical stability of ideal two-dimensional metals in several lattice geometries using first-principles phonon calculations~\cite{ono2020dynamical}.
The study shows that stability depends on the arrangement of the atoms, rather than following from the stability of the bulk element alone.
Hess examined broader periodic trends in two-dimensional bonding among main-group elements~\cite{hess2022periodicity}.
These studies provide structural and bonding guidance, but do not establish that every proposed few-layer configuration is experimentally accessible.
We therefore treat geometry and dynamical stability as questions to be tested for each model.

Chahal et al. reported beryllene sheets obtained by sonochemical exfoliation and characterised several lattice motifs~\cite{chahal2023beryllene}.
These observations provide an experimental basis for studying atomically thin beryllium.
However, an observed sheet and an idealised computational phase are not necessarily the same structure.
The layer number, atomic arrangement, and experimental environment all enter this comparison.
We therefore distinguish the reported experimental geometries from the pristine and hydrogen-functionalized models considered here.
Names such as $\alpha$- and $\beta$-beryllene are used together with a structural definition, since a phase name alone is not sufficient to match models from different studies.

Theoretical calculations have already established a basis for discussing electronic and topological properties in beryllene.
Li et al. investigated a triangular monolayer, denoted as $\alpha$-beryllene, and a corresponding two-layer structure, denoted as $\beta$-beryllene~\cite{li2023coexistence}.
They reported non-trivial band topology in the $\alpha$ phase, while the $\beta$ phase was topologically trivial.
Their electron--phonon calculations also predicted superconducting behaviour in both structures.
These theoretical results provide reference structures and questions against which additional geometries can be examined.
Here, we extend the comparison to a cubic trilayer and selected hydrogen-functionalized configurations.
The question is how changes in layer geometry and surface bonding affect electronic structure and optical response.
For the selected pristine phases, we also examine the interpretation of band topology.

Hydrogen--beryllium interactions have been studied in other structural settings.
Bachurin and Vladimirov investigated hydrogen on beryllium surfaces using \textit{ab initio} calculations~\cite{bachurin2015ab}.
This provides a surface-adsorption reference, but a surface of a bulk crystal retains an underlying three-dimensional environment.
A free-standing layer has different coordination and can respond differently to adsorption.
Separately, Li et al. used global structure search and first-principles calculations to propose two-dimensional beryllium hydride structures~\cite{li2018global}.
Those calculations concern compound formation and changes in the bonding network, rather than the termination of a prescribed elemental beryllium layer.
We therefore distinguish bulk-surface adsorption, hydride-sheet formation, and hydrogen functionalization of a selected parent structure.

Together, these studies motivate a comparison in which both the parent geometry and hydrogen configuration are specified.
They do not determine in advance whether hydrogenation will preserve metallicity, open a gap, or favour a different structure.
Chapter~\ref{cha:project_beryllene} examines these questions for the configurations considered in this thesis.
Structural stability, electronic properties, and optical response are assessed within the corresponding computational approximations.
The band-topology analysis concerns the selected pristine phases.

\subsection{Computational approaches and their limitations}
\label{subsec:lit_methods}

The preceding studies use different calculations to answer different physical questions.
We now compare the information provided by these approaches and the limits that matter for the present work.
The aim here is to identify the basis of a prediction; the detailed theory and calculation settings are given in the later chapters.

Firstly, density functional theory (DFT) provides a practical framework for comparing candidate structures and their electronic states.
Semilocal approximations such as the Perdew--Burke--Ernzerhof (PBE) functional are widely used for these tasks~\cite{perdew1996generalized}.
For example, the phonon calculations of Ono examine how changing the lattice geometry affects the dynamical stability of two-dimensional metals~\cite{ono2020dynamical}.
Such a comparison answers a structural question within a particular approximation.
It does not, by itself, establish the accuracy of an optical excitation energy or a superconducting transition temperature.
We therefore distinguish the property being tested from the general choice to use DFT.

The exchange-correlation treatment is particularly relevant when comparing band gaps and electronic states.
Screened hybrid functionals include a fraction of short-range exact exchange and provide an alternative to semilocal descriptions~\cite{heyd2003hybrid,krukau2006influence}.
Here, we use the 2006 form of the Heyd--Scuseria--Ernzerhof functional, denoted as HSE06.
Comparing PBE and HSE06 in the boron-based studies allows us to assess the sensitivity of the predicted electronic and optical properties to this choice.
However, a hybrid-functional band structure remains an approximate description.
It is not equivalent to a quasiparticle calculation, and an independent-particle optical calculation does not include the attraction between an excited electron and the hole it leaves behind.
Onida et al. distinguish these excitation problems and discuss the roles of many-body self-energy corrections and electron--hole interactions~\cite{onida2002electronic}.
Thygesen further explains how screening and the dielectric environment affect quasiparticles, excitons, and plasmons in two-dimensional materials~\cite{thygesen2017}.
Consequently, agreement between two functionals is useful evidence of a trend, but not a complete validation of an experimental spectrum.

Secondly, a comparison between isolated sheets and heterostructures requires an appropriate treatment of the interface.
Dispersion corrections provide one way to include interactions that are not adequately described by a semilocal functional alone~\cite{grimme2010consistent,grimme2011effect}.
In the graphene-based bilayers considered here, the relaxed separation and atomic registry define the geometry on which the electronic and optical calculations are performed.
Changes in that geometry can therefore affect the inferred role of interfacial coupling.
The optical representation also matters.
Matthes et al. showed how a dielectric tensor obtained from a periodic supercell can be related to the response of an individual sheet or a film of specified thickness, including its out-of-plane response~\cite{matthes2016influence}.
This comparison makes clear that vacuum spacing, effective thickness, and polarization must be stated when reporting optical quantities for two-dimensional models.

Next, we distinguish the different meanings of structural stability.
A total-energy comparison gives an energetic preference relative to the structures or reference states included in that comparison.
A harmonic phonon calculation tests the response to small displacements around a particular geometry~\cite{baroni2001phonons}.
Ab initio molecular dynamics (AIMD) follows the motion of a finite system over a finite time at the chosen simulation conditions.
For example, the evidence used in the prediction of bulk o-\ce{B14} combines several such tests, whereas Ono's survey focuses on harmonic dynamical stability~\cite{han2023superconducting,ono2020dynamical}.
These tests complement one another, but they answer different questions.
They do not replace an assessment of competing phases, formation pathways, or the experimental environment.
We therefore use the terms energetic preference, dynamical stability, and short-time thermal persistence according to the evidence available.

Finally, superconductivity and topology require additional information beyond a ground-state band plot.
The borophene study of Gao et al. estimates transition temperatures from electron--phonon coupling using a McMillan--Allen--Dynes expression, while Li et al. use Migdal--Eliashberg calculations for beryllene~\cite{gao2017prediction,li2023coexistence}.
The resulting temperatures belong to particular materials and treatments of the pairing problem.
They should not be compared without identifying those choices.
Wannier interpolation,
as implemented in EPW, provides a means of evaluating electron--phonon quantities on dense sampling meshes~\cite{ponce2016epw}.
The topology problem instead concerns the structure of electronic wave functions and requires a suitable band subspace and the relevant symmetries~\cite{fu2007topological}.
For the present work, these methods address additional questions in selected systems; they do not serve as interchangeable measures of material performance.

\subsection{Experimental evidence and theoretical models}
\label{subsec:lit_critical_status}

At this stage, the review has introduced both experimentally reported materials and computational models.
We now make the distinction explicit for the systems used in this thesis.

Firstly, graphene, substrate-supported borophene phases, and epitaxial \ce{BC3} sheets have experimental precedents~\cite{novoselov2004electric,zhong2017metastable,chen2022synthesis,yanagisawa2004phonon,tanaka2005novel}.
Borophene--graphene combinations and exfoliated beryllene sheets have also been reported~\cite{liu2019borophene,hou2021borophene,chahal2023beryllene}.
These observations establish relevant material families, while retaining the specific geometries and environments of the experiments.

Secondly, the \ce{B4C3} monolayer used here and the bulk o-\ce{B14} framework have theoretical origins~\cite{tian2019,han2023superconducting}.
The bulk o-\ce{B14} structure is therefore a calculated reference for dimensional reduction, rather than evidence that its monolayer or bilayer has been prepared.
Likewise, the existence of experimental beryllene sheets does not establish the realization of every idealised $\alpha$, $\beta$, cubic trilayer, or hydrogen-functionalized model.

Thirdly, the present calculations concern specified freestanding heterostructures, low-dimensional o-\ce{B14} models, and pristine or hydrogen-functionalized beryllium configurations.
Their structural, electronic, and optical properties are predictions under the stated approximations.
This distinction defines how their contribution is assessed.
A known monolayer can still provide an informative reference for a new interface comparison, and a predicted bulk structure can motivate a new low-dimensional model.
The contribution must identify which structure or property is being investigated and what evidence supports the resulting conclusion.

\section{Motivation for optical and topological analysis}
\label{sec:intro_observables}

% opening paragraph
The previous section discussed the materials and the computational approaches used to study them.
We now explain why electronic structure alone does not answer all the questions considered in this thesis.
Optical response connects electronic transitions to the interaction with light, while band topology concerns the organisation of the electronic states across the Brillouin zone.
These analyses provide different information and are used for different purposes.

\subsection{Optical response as a probe of electronic structure}
\label{subsec:opt_observables_motivation}

The electronic structure describes the states available to the electrons.
However, the band energies alone do not determine the response to light.
We also need to consider which transitions are allowed and how their strength changes with the direction of the electric field.
We therefore compare optical spectra between constituent monolayers and heterostructures, between different dimensional forms of o-\ce{B14}, and between pristine and hydrogen-functionalized beryllium structures.
These comparisons test whether structural changes produce distinguishable changes in the calculated response.

Firstly, we compare transition energies and spectral weight.
Within the adopted independent-particle approximation, these quantities describe the electronic response at different photon energies~\cite{onida2002electronic}.
Comparing an isolated layer with a bilayer helps identify changes associated with interfacial coupling.
Comparing pristine and hydrogen-functionalized structures instead examines how changes in surface bonding are reflected in the spectrum.
The purpose is to connect a change in the atomic model to a change in the optical response.

We then compare different polarizations to examine optical anisotropy.
The dielectric response within the structural plane can differ from the response perpendicular to it, and inequivalent in-plane directions may also need to be distinguished.
This directional comparison connects the optical calculation to the symmetry and geometry of the model.
Off-diagonal tensor components must be interpreted in the chosen coordinate system and with the relevant symmetry constraints.
Their presence alone does not establish optical activity.
The derived absorption coefficient, refractive index, extinction coefficient, and reflectivity describe further aspects of the same dielectric response.
They therefore provide related descriptions of light propagation through the same dielectric response.

Next, we examine the energy-loss features.
Here, a loss peak can indicate a collective excitation, but its interpretation also depends on dielectric screening and damping.
We therefore consider the loss spectrum together with the real and imaginary dielectric response~\cite{thygesen2017}.
A peak or a zero crossing alone does not establish a particular plasmon mode.
For two-dimensional systems, the relation to a measured loss spectrum also depends on momentum transfer, geometry, and the surrounding environment.
These conditions define how the calculated features can be compared with electron energy-loss spectroscopy (EELS).

Finally, we specify the limits of the optical comparison.
The calculated spectra depend on the treatment of intraband transitions and electron--hole interactions~\cite{onida2002electronic}.
For periodic two-dimensional models, they also depend on the convention used for vacuum and effective thickness~\cite{matthes2016influence}.
These choices determine how the calculated response can be related to a measurement in a particular environment.
Chapter~\ref{cha:fundamentals_b} develops the dielectric response and its relation to optical properties.
The research chapters state the computational treatment used for their respective systems.

\subsection{Physical significance of band topology}
\label{subsec:topo_motivation}

Optical spectra describe the response of the electronic system.
We now consider another question: how are the electronic states connected across the Brillouin zone?
Band topology addresses this question through properties of the wave functions that are not contained in the size of a band gap alone.
For a two-dimensional insulating system with time-reversal symmetry and the appropriate spinful band description, a non-trivial $\mathbb{Z}_2$ invariant is associated with helical edge states~\cite{fu2007topological}.
This connection motivates the investigation of topology in light-element structures, where structural and orbital changes may be relevant even though atomic spin--orbit coupling is weak.

However, the conditions of the classification must be stated.
In an inversion-symmetric system, the Fu--Kane criterion evaluates the invariant using the parity eigenvalues of Kramers pairs at the time-reversal-invariant momenta~\cite{fu2007topological}.
This provides an efficient evaluation once the relevant band subspace has been defined.
For a metallic system, its separation from the other bands must hold through a direct gap throughout the Brillouin zone.
Such a subspace can exist even when indirect band overlap makes the material metallic.
Its topology does not by itself establish quantum spin Hall insulating transport in the material.
Chapter~\ref{cha:fundamentals_c} develops these symmetry and band-isolation conditions.
They provide the basis for assessing the topological interpretation of the beryllium structures considered in chapter~\ref{cha:project_beryllene}.

\section{Research questions and scope of the thesis}
\label{sec:intro_gaps}

% opening paragraph
The preceding discussion identifies a common question.
How do changes in atomic structure modify the electronic structure and optical response of light-element two-dimensional materials?
We investigate this question using first-principles calculations.

The three projects examine this question in different bonding environments.
Each project compares reference structures with modified structures.
The modification may involve an interface, a change in dimensionality, or hydrogen termination.
This shared comparison connects the projects.
Magnetism, superconductivity, and band topology enter as further questions for selected systems.

More specifically, we address the following research questions:

\begin{enumerate}
\item The modification of electronic structure and polarization-dependent optical response induced by coupling graphene to selected boron-based monolayers,
evaluated relative to the isolated constituents.
\item The influence of dimensional reduction and hydrogen termination on electronic, magnetic,
and electron--phonon properties within the pentagonal-bipyramid boron framework, along with the corresponding signatures in the optical response.
\item The effects of layer geometry and hydrogen functionalization on the stability, electronic structure,
and optical anisotropy of the examined beryllium configurations, together with the topological character of the selected pristine phases.
\end{enumerate}

The first question builds on known monolayers and experimental precedents for related interfaces.
Its contribution is the comparison of the selected graphene-based bilayers and their optical response within a stated computational framework.
The second question uses the previously predicted bulk o-\ce{B14} structure as a reference and examines its low-dimensional and hydrogen-terminated forms.
The third question extends the comparison of beryllium structures to a cubic trilayer and selected hydrogen-functionalized configurations, while retaining earlier $\alpha$- and $\beta$-beryllene studies as references.
Together, these comparisons identify which changes in response follow the structural modifications studied and where the interpretation depends on the method or model.

For the present work, we compare selected structural models using specified computational approximations.
The calculations provide evidence for the properties of those models.
We therefore state the limits of each prediction and distinguish computational results from experimental observations.
Structural relaxation, phonons, and finite-time dynamics are used to assess different aspects of stability, rather than to imply experimental synthesizability.
Similarly, the magnetic, superconducting, and topological analyses are interpreted within the corresponding electronic states and theoretical descriptions.

\newpage
\section{Thesis organisation}
\label{sec:intro_organisation}

% opening paragraph
The discussion follows the sequence of these research questions.
We first introduce the theoretical framework and then turn to the three material studies.

Chapter~\ref{cha:fundamentals_a} introduces the many-body problem and the density-functional framework used in the calculations.
It also discusses the exchange-correlation approximations used later.
Chapter~\ref{cha:fundamentals_b} then develops the linear optical response.
It connects the dielectric response to the optical properties examined in the research chapters.
Chapter~\ref{cha:fundamentals_c} develops the band-subspace and symmetry conditions for the topological analysis.
It distinguishes the topology of an isolated band subspace from an insulating state at the Fermi level.

In chapter~\ref{cha:project_bilayers}, we compare graphene-based heterostructures with their constituent monolayers.
This comparison examines how interfacial coupling changes the electronic structure and optical response.

We next consider the pentagonal-bipyramid boron framework in chapter~\ref{cha:project_boron}.
Here, we examine the effects of dimensional reduction and hydrogen termination on electronic, magnetic, electron--phonon, and optical properties.

Chapter~\ref{cha:project_beryllene} extends the discussion to pristine and hydrogen-functionalized beryllium structures, including a cubic trilayer.
We examine their structural stability, electronic properties, and optical response, together with band topology in the selected pristine phases.

Finally, chapter~\ref{cha:conclusion_outlook} brings the findings together.
It outlines further work suggested by these results.
